% ============================================================
% MainPaper.tex — Artículo IEEE, Etapa 3 (final)
% IC-6200 Inteligencia Artificial · II Semestre 2026
% Build: latexmk -pdf MainPaper.tex
% ============================================================
\documentclass[conference]{IEEEtran}

% ============================================================
% preamble.tex — configuración compartida, compatible con IEEEtran
%
% NO agregar: geometry, fancyhdr, setspace, tocloft, natbib,
% caption, subcaption, cleveref. IEEEtran los maneja internamente.
% ============================================================

% ---------- Idioma ----------
\usepackage[T1]{fontenc}
\usepackage[spanish,es-tabla,es-noquoting,es-nodecimaldot]{babel}
\usepackage{microtype}

% ---------- Encabezados de IEEEtran en español ----------
\renewcommand{\abstractname}{Resumen}
\renewcommand{\IEEEkeywordsname}{Palabras clave}
\renewcommand{\refname}{Referencias}
\renewcommand{\figurename}{Fig.}
\renewcommand{\tablename}{TABLA}

% ---------- Numeración de página ----------
% IEEEtran conference la desactiva por defecto porque en actas
% reales la asigna el editor. Se activa por ser trabajo de curso.
\pagestyle{plain}
\AtBeginDocument{\thispagestyle{plain}}

\usepackage{amsmath}
\usepackage{amssymb}
\usepackage{amsthm}
\usepackage{mathtools}
\usepackage{siunitx}
\sisetup{per-mode=symbol}

\usepackage{cite}

\usepackage{xcolor}
\definecolor{tecnavy}{RGB}{31,56,100}
\definecolor{tecsteel}{RGB}{46,85,150}
\definecolor{tecrule}{RGB}{155,184,216}
\definecolor{boxbg}{RGB}{242,246,252}
\definecolor{labelgray}{RGB}{120,120,120}
\definecolor{notaverde}{RGB}{22,110,60}
\definecolor{decisionambar}{RGB}{170,110,10}
\definecolor{codebg}{RGB}{248,248,248}
\definecolor{codekw}{RGB}{0,0,180}
\definecolor{codecmt}{RGB}{0,120,60}
\definecolor{codestr}{RGB}{170,55,55}

\usepackage{graphicx}
\graphicspath{{img/}{figures/}{./}}
\usepackage{tikz}
\usepackage{pgfplots}
\pgfplotsset{compat=1.18}
\usetikzlibrary{arrows.meta,positioning,shapes.geometric,calc,fit,
                backgrounds,decorations.pathreplacing}

\usepackage{array}
\usepackage{booktabs}
\usepackage{multirow}
\usepackage{makecell}
\usepackage{tabularx}
\usepackage{threeparttable}
\usepackage{enumitem}

\newcolumntype{Y}{>{\raggedright\arraybackslash}X}
\newcolumntype{P}[1]{>{\raggedright\arraybackslash}p{#1}}
\renewcommand{\arraystretch}{1.15}
\setlist[itemize]{leftmargin=*,noitemsep,topsep=2pt}
\setlist[enumerate]{leftmargin=*,noitemsep,topsep=2pt}

\theoremstyle{definition}
\newtheorem{definicion}{Definición}
\newtheorem{ejemplo}{Ejemplo}
\theoremstyle{plain}
\newtheorem{proposicion}{Proposición}
\theoremstyle{remark}
\newtheorem*{observacion}{Observación}

\usepackage{listings}
\lstdefinestyle{code}{
  basicstyle=\ttfamily\scriptsize,
  keywordstyle=\color{codekw}\bfseries,
  commentstyle=\color{codecmt}\itshape,
  stringstyle=\color{codestr},
  backgroundcolor=\color{codebg},
  numbers=left,
  numberstyle=\tiny\color{labelgray},
  numbersep=6pt,
  frame=single,
  rulecolor=\color{tecrule},
  breaklines=true,
  postbreak=\mbox{\textcolor{tecsteel}{$\hookrightarrow$}\space},
  showstringspaces=false,
  tabsize=2,
  captionpos=b,
  columns=flexible,
  upquote=true
}
\lstset{style=code}

\usepackage[hidelinks]{hyperref}

% ---------- Entrada segura ----------
\newcommand{\cajaarchivofaltante}[1]{%
  \par\medskip\noindent
  \fcolorbox{red!60!black}{red!5}{%
    \parbox{\dimexpr\linewidth-2\fboxsep-2\fboxrule}{%
      \centering\ttfamily\scriptsize ARCHIVO FALTANTE: \detokenize{#1}}}%
  \par\medskip}

\newcommand{\cajaimagenfaltante}[1]{%
  \begingroup\setlength{\fboxsep}{0pt}%
  \fcolorbox{tecrule}{boxbg}{%
    \parbox[c][2.4cm][c]{0.9\linewidth}{\centering\ttfamily\scriptsize
      IMAGEN NO ENCONTRADA\\[3pt]\detokenize{#1}}}%
  \endgroup}

\newcommand{\safeinput}[1]{%
  \IfFileExists{#1.tex}{\input{#1}}{%
    \PackageWarning{tecpaper}{Falta el archivo de sección #1.tex}%
    \cajaarchivofaltante{#1.tex}}}

\newcommand{\safeimage}[2][width=\columnwidth]{%
  \IfFileExists{#2}{\includegraphics[#1]{#2}}{%
    \IfFileExists{img/#2}{\includegraphics[#1]{img/#2}}{%
      \PackageWarning{tecpaper}{Falta la imagen #2}%
      \cajaimagenfaltante{#2}}}}

\newcommand{\code}[1]{\texttt{#1}}

% ============================================================
% MARCADORES DE REDACCIÓN
%
% Los tres se apagan antes de entregar redefiniéndolos como vacíos:
%   \renewcommand{\pc}[1]{}
%   \renewcommand{\razon}[1]{}
%   \renewcommand{\decidir}[1]{}
% Recompilar: todo lo que desaparezca del PDF nunca se escribió.
% ============================================================

% \pc  = POR COMPLETAR. Contenido que depende de los ejes A y C.
\newcommand{\pc}[1]{%
  \texorpdfstring{\textcolor{red}{\textbf{[POR COMPLETAR: #1]}}}%
                 {[POR COMPLETAR: #1]}}

% \razon = POR QUÉ esta parte debe decir lo que dice, con su fuente.
\newcommand{\razon}[1]{%
  \texorpdfstring{\textcolor{notaverde}{\small\textit{[RAZÓN: #1]}}}%
                 {[RAZÓN: #1]}}

% \decidir = requiere decisión de equipo o criterio humano.
\newcommand{\decidir}[1]{%
  \texorpdfstring{\textcolor{decisionambar}{\small\textbf{[DECIDIR: #1]}}}%
                 {[DECIDIR: #1]}}

% ---------- Columnas X con peso, para tabularx ----------
% Permite repartir el ancho flexible en proporciones distintas.
% Los pesos de todas las columnas W deben sumar el número de
% columnas W de la tabla.
\newcolumntype{W}[1]{>{\hsize=#1\hsize\raggedright\arraybackslash}X}

% ---------- Marcador compacto para celdas de tabla ----------
% \pc dentro de una celda produce una cadena larga que rompe el
% ajuste de columnas. Este marcador ocupa tres caracteres.
% Se apaga igual que los demás: \renewcommand{\pct}{}
\newcommand{\pct}{\textcolor{red}{\textbf{[?]}}}

% ============================================================
% APAGADO DE MARCADORES
% Descomentar las cuatro líneas antes de entregar y recompilar.
% Todo lo que desaparezca del PDF nunca se escribió.
% ============================================================

%\renewcommand{\pc}[1]{}
%\renewcommand{\razon}[1]{}
%\renewcommand{\decidir}[1]{}
%\renewcommand{\pct}{}

\begin{document}

\title{Estimación de la Probabilidad de Victoria a partir del Estado de
Juego en el Pok\'emon Trading Card Game: un Estudio de Clasificación
Tabular con Control Explícito de Fuga de Información\\
{\footnotesize \textit{Etapa 3: modelo, agente y validación ---
Curso IC-6200 Inteligencia Artificial, II Semestre 2026}}}

\author{
\IEEEauthorblockN{Adriel S. Chaves Salazar}
\IEEEauthorblockA{\textit{Escuela de Ingenier\'ia en Computadores} \\
\textit{Instituto Tecnol\'ogico de Costa Rica}\\
Cartago, Costa Rica \\
2021031465}
\and
\IEEEauthorblockN{Daniel Duarte Cordero}
\IEEEauthorblockA{\textit{Escuela de Ingenier\'ia en Computadores} \\
\textit{Instituto Tecnol\'ogico de Costa Rica}\\
Cartago, Costa Rica \\
2022012866}
\and
\IEEEauthorblockN{Sebasti\'an Hern\'andez Bonilla}
\IEEEauthorblockA{\textit{Escuela de Ingenier\'ia en Computadores} \\
\textit{Instituto Tecnol\'ogico de Costa Rica}\\
Cartago, Costa Rica \\
2022093651}
}

\maketitle

\begin{abstract}
El Pok\'emon Trading Card Game es un juego de cartas por turnos en el que
cada jugador observa el tablero completo y su propia mano, pero no la mano
del oponente ni el contenido de las cartas de premio. Determinar quién lleva
ventaja exige considerar los premios restantes, los puntos de vida, las
energías y los estados alterados, lo que puede ser difícil para un observador
sin experiencia. Este trabajo utiliza aprendizaje supervisado para estimar
la probabilidad de victoria del jugador que actúa a partir de un estado
observable. Para ello emplea un corpus
público de repeticiones de partidas entre agentes automáticos. La dificultad
principal está en la evaluación. Una partida produce cientos de estados
sucesivos que difieren en una sola acción y comparten el mismo resultado. Si
estos estados se reparten al azar entre
entrenamiento y prueba, el modelo se evaluaría sobre situaciones
casi idénticas a las que ya vio. Por esa razón, se emplea una
representación tabular con una fila por estado y 62 características
restringidas a la información pública de ese instante, una comparación de
modelos clásicos desde un predictor trivial hasta potenciación del
gradiente, y particiones que mantienen cada partida completa dentro de un
único subconjunto. Se evalúan tres escenarios que comprenden partidas nuevas,
arquetipos de mazo no vistos y periodos posteriores. Sobre 55\,196 partidas, el modelo de
potenciación del gradiente alcanza en prueba un AUC de 0.833 sobre partidas
nuevas y de 0.789 y 0.787 bajo cambio temporal y arquetipos no vistos, y
supera a todas las referencias con intervalos que excluyen el cero. La
capacidad predictiva crece de 0.68 al inicio a 0.94 en la etapa final. En
mazos conocidos, al comienzo de la partida el modelo no predice mejor que una
referencia que solo considera quién juega primero y qué Pok\'emon están
activos. En mazos no vistos durante el entrenamiento, sí aporta información
adicional. Bajo cambio temporal, contra lo esperado, al modelo le cuesta más
distinguir los estados favorables de los desfavorables que asignar
probabilidades acordes con los resultados reales. El análisis SHAP muestra
que el modelo evalúa la posición según el balance de recursos, como los
premios, los descartes y el mazo.
Finalmente, el modelo se integra en un agente con cuatro roles
(planificador, ejecutor, verificador y redactor) en el que toda cifra
sale del modelo y pasa por un verificador determinista. En una validación
con 35 casos automáticos y 31 consultas al agente completo, todas las
respuestas finales tuvieron cifras trazables. El límite principal es que
la mayoría de los errores muy seguros de las etapas media y final llegan
al usuario sin advertencia.
\end{abstract}

\begin{IEEEkeywords}
estimación de probabilidad de victoria, juegos de cartas coleccionables,
aprendizaje automático tabular, fuga de datos, validación cruzada agrupada,
calibración de clasificadores, agentes con herramientas, machine learning
\end{IEEEkeywords}

\safeinput{sec/01introduccion}
\safeinput{sec/02trabajos_relacionados}
\safeinput{sec/03baseline_literatura}
\safeinput{sec/04metodologia}
\safeinput{sec/05analisis_tecnico}
\safeinput{sec/07resultados}
\safeinput{sec/10agente}
\safeinput{sec/08discusion}
\safeinput{sec/06etica}
\safeinput{sec/09conclusiones}

\bibliographystyle{IEEEtran}
\bibliography{sec/99referencias}

\end{document}
